\section{Antecedentes del proyecto}

\subsection{Situación previa}

Antes del inicio de la implementación, la realidad operativa de \textbf{Machu Picchu Exclusive Tours E.I.R.L.} se caracterizaba por:

\begin{itemize}
    \item Un \textbf{flujo de ventas completamente virtual}, apoyado en conversaciones de \textit{WhatsApp}, videollamadas y recomendaciones en eventos, sin un sistema de registro centralizado.
    \item \textbf{Gestión de datos sensibles} (principalmente fotos de pasaportes de los PAX) recibidos por mensajería instantánea y borrados periódicamente (cada uno o dos años) por motivos de seguridad, lo que impedía construir una base de datos estructurada.
    \item Ausencia de un \textbf{historial consolidado} de las interacciones con los turistas (notas, acuerdos, cambios de itinerario), lo que dificultaba el seguimiento interno y la estandarización del servicio.
    \item Manejo flexible de \textbf{adelantos y pagos en cuotas} (dos o tres cuotas), con riesgo de perder visibilidad sobre saldos pendientes antes de ejecutar pagos no reembolsables (boletos de tren, entradas a Machu Picchu, hoteles).
\end{itemize}

\subsection{Problemática identificada}

El diagnóstico situacional permitió identificar las siguientes problemáticas críticas:

\begin{itemize}
    \item \textbf{Fragmentación de la información}: Datos de pasaportes, itinerarios, adelantos y confirmaciones dispersos en múltiples chats, sin un repositorio único y estructurado.
    \item \textbf{Riesgo de errores en datos sensibles}: Cualquier error en nombres, fechas o números de pasaporte al momento de comprar boletos o entradas a Machu Picchu es difícil de corregir y puede afectar gravemente la experiencia del turista.
    \item \textbf{Control limitado de adelantos y saldos}: La agencia asume un riesgo financiero al realizar compras no reembolsables si no se verifica de manera rigurosa el estado de pagos del cliente.
    \item \textbf{Ausencia de visión 360° del PAX}: Sin un CRM turístico, se dificulta tener una imagen completa del historial del turista (viajes realizados, preferencias, fechas relevantes, encuestas post--viaje).
    \item \textbf{Dificultad para escalar la operación}: La dependencia de la memoria y del seguimiento manual limita la capacidad de crecimiento ordenado y la delegación efectiva de funciones dentro del equipo.
\end{itemize}

